\section{稳态分布映射与不变测度条件}

本阶段将稳态分布模块改写为完全显式的测度论形式：对象先定义后使用，定理调用前先列前提，关键不等式逐行展开。

\subsection{Benchmark Fidelity Check}

\paragraph{Step 0：先写出 benchmark 的稳态测度方程与定义域。}
主文稿把稳态分布的运动方程写成
\begin{equation}
\mu'_j(B)
=
\sum_{s\in\{A,B,C\}}\int_Z \mathbf 1\{g_s(z;M)=j\}\,Q_j(B\mid z,M)\,\mu_s(dz)
+n_E^j(M)G(B),
\qquad j\in\{A,B,C\},
\label{eq:p7zh_main_stationary_measure}
\end{equation}
其中 $B\subseteq Z$ 是任意可测集，$g_s(z;M)$ 在主文稿中被写成 optimal organizational policy。全文固定制度与价格：
\[
M\in\{0,1\},
\qquad
(w,p_m)\in\mathcal W\times\mathcal P_m,
\]
其中
\[
\mathcal P_m=[\underline p_m,\overline p_m]\subset\mathbb R_+,
\qquad
\mathcal W=[\underline w,\overline w]\subset\mathbb R_+,
\]
并满足 $0\le\underline p_m\le\overline p_m<\infty$、$0\le\underline w\le\overline w<\infty$。

状态空间是紧度量空间 $(Z,\mathcal B(Z))$，活跃组织类型为 $\{A,B,C\}$。

\paragraph{Step 1：精确对象与 policy representation 的张力。}
Phase 4 保留的 benchmark choice-region 写法是
\[
\mathcal Z_{sj}(M)=\{z\in Z: j \text{ solves the Bellman problem for current type }s\}.
\]
但 \eqref{eq:p7zh_main_stationary_measure} 使用的是单值 policy rule $g_s$。只有在 Bellman maximizer 唯一，或者另外加入可测 tie-breaking convention 时，这两种表示才会自动一致。因此，本节以下推导是条件于 selector representation 的写法，它对应主文稿中的稳态测度方程，但不应被误读为“Phase 4 的 set-based choice regions 自动变成了 partition”。

\paragraph{Step 2：本节使用的继承对象。}
\begin{enumerate}
    \item Phase 3 的 argmax correspondence
    \[
    \Gamma_s(z;M):=\arg\max_{j\in\{A,B,C,\mathrm{exit}\}}W_{sj}(z;M),
    \qquad s\in\{A,B,C\};
    \]
    \item 一个可测 selector representation
    \[
    g_s:Z\to\{A,B,C,\mathrm{exit}\},
    \qquad s\in\{A,B,C\},
    \]
    满足对每个 $z$ 都有 $g_s(z;M)\in\Gamma_s(z;M)$；
    \item 转移核 $Q_j(\cdot\mid z,M)$，其中 $j\in\{A,B,C\}$；
    \item entrant masses $n_E^A(M),n_E^B(M),n_E^C(M)$ 与 entrant distribution $G$。
\end{enumerate}

\paragraph{Step 3：benchmark-fidelity 边界。}
下面的稳态测度映射忠实对应主文稿中的 selector-based notation。它不是关于 Phase 4 重叠 choice regions 的自动 row-stochasticity 定理；只要 ties 在经济上相关的状态集合上出现，selector interpretation 就属于补充 convention。

\paragraph{Step 4：由组织策略诱导继续经营核。}
对每个当前类型 $s\in\{A,B,C\}$ 与下一期活跃类型 $j\in\{A,B,C\}$，定义
\begin{equation}
K_{sj}(B\mid z,M):=\mathbf 1\{g_s(z;M)=j\}\,Q_j(B\mid z,M),
\qquad B\in\mathcal B(Z).
\label{eq:p7zh_Ksj}
\end{equation}

\paragraph{Step 2：定义测度空间与动态映射。}
记 $\mathcal M_+(Z)$ 为 $Z$ 上有限非负 Borel 测度集合，$\mathcal M(Z)$ 为有限符号 Borel 测度集合。
定义
\[
\mathcal M_+^3:=\mathcal M_+(Z)^3,
\qquad
\mathcal M^3:=\mathcal M(Z)^3.
\]

对任意 $\mu=(\mu_A,\mu_B,\mu_C)\in\mathcal M_+^3$，定义
\begin{equation}
(\mathcal T_M\mu)_j(B):=\sum_{s\in\{A,B,C\}}\int_Z K_{sj}(B\mid z,M)\,\mu_s(dz),
\qquad j\in\{A,B,C\}.
\label{eq:p7zh_T}
\end{equation}
定义进入流入测度向量
\begin{equation}
\nu_j(B):=n_E^j(M)\,G(B),
\qquad j\in\{A,B,C\}.
\label{eq:p7zh_nu}
\end{equation}
定义一步映射
\begin{equation}
\Phi_M(\mu):=\mathcal T_M\mu+\nu.
\label{eq:p7zh_Phi}
\end{equation}
稳态分布向量定义为不动点
\begin{equation}
\mu^*=\Phi_M(\mu^*).
\label{eq:p7zh_fixed_point}
\end{equation}

\subsection{假设}

\begin{assumption}[SD0：可测 selector representation]
对每个当前类型 $s\in\{A,B,C\}$，存在可测映射
\[
g_s:Z\to\{A,B,C,\mathrm{exit}\}
\]
使得对每个 $z\in Z$ 都有 $g_s(z;M)\in\Gamma_s(z;M)$。如果 Bellman maximizer 唯一，则 SD0 自动成立；如果存在 ties，则 SD0 是补充性的 measurable-selection convention。
\end{assumption}

\begin{assumption}[SD1：核可测]
对每个 $s,j\in\{A,B,C\}$ 与每个 $B\in\mathcal B(Z)$，映射 $z\mapsto K_{sj}(B\mid z,M)$ 是 Borel 可测。
\end{assumption}

\begin{assumption}[SD2：进入流有限]
$n_E^A(M),n_E^B(M),n_E^C(M)$ 有限且非负，$G$ 为概率测度。
\end{assumption}

\begin{assumption}[SD3：算子层面的压缩正则条件]
由 \eqref{eq:p7zh_T} 定义的 $\mathcal T_M$ 在符号测度向量空间上的线性延拓满足
\begin{equation}
\|\mathcal T_M\eta\|_{\infty,\mathrm{TV}}
\le
\bar\rho\,\|\eta\|_{\infty,\mathrm{TV}}
\qquad\text{对每个 }\eta\in\mathcal M^3,
\label{eq:p7zh_operator_contraction}
\end{equation}
其中某个常数 $\bar\rho\in[0,1)$。这是关于 selector-induced stationary-distribution operator 的补充正则条件；它并不会由 benchmark 中“每个 $Q_j(\cdot\mid z,M)$ 都是 probability kernel”这一事实自动推出。
\end{assumption}

\subsection{范数与完备性}

对符号测度向量 $\eta=(\eta_A,\eta_B,\eta_C)\in\mathcal M^3$，定义 max-TV 范数
\begin{equation}
\|\eta\|_{\infty,\mathrm{TV}}:=\max_{s\in\{A,B,C\}}\|\eta_s\|_{\mathrm{TV}},
\qquad
\|\eta_s\|_{\mathrm{TV}}:=|\eta_s|(Z).
\label{eq:p7zh_norm}
\end{equation}

\begin{lemma}[完备性]
$(\mathcal M^3,\|\cdot\|_{\infty,\mathrm{TV}})$ 是 Banach 空间，且 $\mathcal M_+^3$ 在其中为闭子集。
\end{lemma}

\begin{proof}
每个 $\mathcal M(Z)$ 在 TV 范数下是 Banach。有限笛卡尔积在 max 范数下仍是 Banach。非负性在 TV 极限下保持，所以 $\mathcal M_+^3$ 是闭子集。
\end{proof}

\subsection{映射 $\Phi_M$ 的良定义}

\begin{lemma}[良定义]
在 SD0--SD2 下，若 $\mu\in\mathcal M_+^3$，则 $\Phi_M(\mu)\in\mathcal M_+^3$。
\end{lemma}

\begin{proof}
固定 $j$ 与 $B\in\mathcal B(Z)$。由 SD1，\eqref{eq:p7zh_T} 的被积函数可测。且
\[
0\le K_{sj}(B\mid z,M)\le1,
\]
所以
\[
0\le\int_Z K_{sj}(B\mid z,M)\,\mu_s(dz)\le\mu_s(Z)<\infty.
\]
对 $s$ 求和后，$(\mathcal T_M\mu)_j(B)$ 有限且非负。由核积分标准结论，$(\mathcal T_M\mu)_j$ 关于 $B$ 可列可加，故为有限非负测度。

由 SD2，$\nu_j(B)=n_E^j(M)G(B)$ 是有限非负测度，因此
\[
(\Phi_M(\mu))_j=(\mathcal T_M\mu)_j+\nu_j
\]
仍是有限非负测度。对每个 $j$ 成立，故 $\Phi_M(\mu)\in\mathcal M_+^3$。
\end{proof}

\subsection{压缩不等式（完整代数）}

\begin{lemma}[在位算子的压缩性]
在 SD0--SD3 下，对任意 $\mu,\tilde\mu\in\mathcal M_+^3$，
\begin{equation}
\|\mathcal T_M\mu-\mathcal T_M\tilde\mu\|_{\infty,\mathrm{TV}}
\le\bar\rho\,\|\mu-\tilde\mu\|_{\infty,\mathrm{TV}}.
\label{eq:p7zh_contraction_T}
\end{equation}
进而
\begin{equation}
\|\Phi_M(\mu)-\Phi_M(\tilde\mu)\|_{\infty,\mathrm{TV}}
\le\bar\rho\,\|\mu-\tilde\mu\|_{\infty,\mathrm{TV}}.
\label{eq:p7zh_contraction_Phi}
\end{equation}
\end{lemma}

\begin{proof}
记
\[
\delta:=\mu-\tilde\mu\in\mathcal M^3.
\]
由于每个 $K_{sj}(B\mid\cdot,M)$ 都是有界可测函数，\eqref{eq:p7zh_T} 中的积分可以通过 Jordan decomposition 从非负测度延拓到符号测度，因此 $\mathcal T_M\delta$ 是良定义的。又因为积分对测度是线性的，
\[
\mathcal T_M\mu-\mathcal T_M\tilde\mu
=
\mathcal T_M(\mu-\tilde\mu)
=
\mathcal T_M\delta.
\]
把 SD3 应用于 $\eta=\delta$，得到
\[
\|\mathcal T_M\mu-\mathcal T_M\tilde\mu\|_{\infty,\mathrm{TV}}
=
\|\mathcal T_M\delta\|_{\infty,\mathrm{TV}}
\le
\bar\rho\,\|\delta\|_{\infty,\mathrm{TV}}
=
\bar\rho\,\|\mu-\tilde\mu\|_{\infty,\mathrm{TV}}.
\]
这正是 \eqref{eq:p7zh_contraction_T}。

最后
\[
\Phi_M(\mu)-\Phi_M(\tilde\mu)
=(\mathcal T_M\mu+\nu)-(\mathcal T_M\tilde\mu+\nu)
=\mathcal T_M\mu-\mathcal T_M\tilde\mu,
\]
故 \eqref{eq:p7zh_contraction_Phi} 成立。
\end{proof}

\subsection{不变测度的存在唯一性}

\begin{proposition}[唯一不变测度向量]
在 SD0--SD3 下，存在唯一 $\mu^*\in\mathcal M_+^3$ 满足
\[
\mu^*=\Phi_M(\mu^*).
\]
\end{proposition}

\begin{proof}
调用 Banach 不动点定理前，逐条核对：
\begin{enumerate}
    \item 完备度量空间：由完备性引理，$\mathcal M_+^3$ 在 $\|\cdot\|_{\infty,\mathrm{TV}}$ 下完备；
    \item 自映射：由“良定义”引理，$\Phi_M:\mathcal M_+^3\to\mathcal M_+^3$；
    \item 压缩：由 \eqref{eq:p7zh_contraction_Phi}，压缩模数为 $\bar\rho<1$。
\end{enumerate}
因此 Banach 定理给出唯一不动点 $\mu^*$。
\end{proof}

\begin{proposition}[总质量上界]
设 $\mu^*$ 为唯一不动点，则
\begin{equation}
\|\mu^*\|_{\infty,\mathrm{TV}}
\le\frac{\|\nu\|_{\infty,\mathrm{TV}}}{1-\bar\rho}.
\label{eq:p7zh_mass_bound}
\end{equation}
\end{proposition}

\begin{proof}
由不动点方程与三角不等式：
\begin{align*}
\|\mu^*\|_{\infty,\mathrm{TV}}
&=\|\mathcal T_M\mu^*+\nu\|_{\infty,\mathrm{TV}}\\
&\le\|\mathcal T_M\mu^*\|_{\infty,\mathrm{TV}}+\|\nu\|_{\infty,\mathrm{TV}}\\
&\le\bar\rho\|\mu^*\|_{\infty,\mathrm{TV}}+\|\nu\|_{\infty,\mathrm{TV}}.
\end{align*}
移项并除以 $1-\bar\rho>0$，得 \eqref{eq:p7zh_mass_bound}。
\end{proof}

\subsection{与转移矩阵对象的接口}

\paragraph{benchmark choice regions 与 selector regions 的区分。}
Phase 4 中的 benchmark choice regions 写成
\[
\mathcal Z_{sj}(M):=\{z\in Z:j\in\Gamma_s(z;M)\},
\qquad s,j\in\{A,B,C,\mathrm{exit}\}.
\]
在 ties 存在时，这些 benchmark regions 可能重叠。本节的 selector representation 则产生辅助的 selector regions
\[
\widehat Z_{sj}(M):=\{z\in Z:g_s(z;M)=j\},
\qquad s,j\in\{A,B,C,\mathrm{exit}\}.
\]
若 Bellman maximizer 唯一，则两者一致；若存在 ties，则 $\widehat Z_{sj}(M)$ 取决于 SD0 中的补充 selector convention。

\paragraph{selector-based transition shares。}
对任何满足 $\mu_s^*(Z)>0$ 的当前类型 $s$，定义辅助的 selector-based transition share
\begin{equation}
\widehat P_{sj}(M):=\frac{\mu_s^*(\widehat Z_{sj}(M))}{\mu_s^*(Z)}.
\label{eq:p7zh_Phat}
\end{equation}

\begin{lemma}[selector-based shares 的行和]
若 $\mu_s^*(Z)>0$，则
\[
\sum_{j\in\{A,B,C,\mathrm{exit}\}}\widehat P_{sj}(M)=1.
\]
\end{lemma}

\begin{proof}
由于 $g_s$ 是单值函数，集合族 $\{\widehat Z_{sj}(M):j\in\{A,B,C,\mathrm{exit}\}\}$ 构成 $Z$ 的可测分割：
\[
Z=\bigsqcup_j \widehat Z_{sj}(M).
\]
由互不相交并的有限可加性，
\[
\mu_s^*(Z)=\sum_j\mu_s^*(\widehat Z_{sj}(M)).
\]
再除以正数 $\mu_s^*(Z)$，得到
\[
1=\sum_j\frac{\mu_s^*(\widehat Z_{sj}(M))}{\mu_s^*(Z)}
=\sum_j \widehat P_{sj}(M).
\]
\end{proof}

\paragraph{如何映回 benchmark object。}
如果对每个当前类型 $s$，Bellman maximizer 在 $\mu_s^*$-几乎处处唯一，则
\[
\widehat Z_{sj}(M)=\mathcal Z_{sj}(M)
\quad\text{up to }\mu_s^*\text{-null sets},
\]
因此 \eqref{eq:p7zh_Phat} 中的 selector-based shares 与 Phase 4 的 set-based benchmark transition shares 一致。若没有这个额外条件，则这里的行和为 1 结论只是辅助 selector representation 的结论，而不是关于重叠 benchmark choice regions 的直接定理。

\subsection{What did we just do, and why does it matter?}

本阶段把“稳态分布”这一直观概念改写成一个完全显式的 measure-valued operator 问题。核心对象是映射 $\Phi_M$：给定候选 incumbent distribution、selector-based policy representation 与 entrant inflows，它返回下一期的 incumbent distribution。一旦这个算子被写清楚，稳态分布问题就变成固定点问题
\[
\mu^*=\Phi_M(\mu^*).
\]
这里把压缩不等式逐行展开并不是形式主义。这样做的作用是明确说明：后续均衡阶段要调用的稳态对象，不只是“存在某个分布”，而是一个在给定条件下单值、稳定、可以通过迭代恢复的对象。
同样重要的是，本版文字现在明确写出唯一性究竟来自哪里：它不是从 benchmark continuation kernels 只是 probability measures 这一点自动推出的，而是来自 SD3 中额外声明的算子层面压缩正则条件。

同样关键的是，本阶段没有把 benchmark choice regions 与 auxiliary selector regions 混为一谈。Phase 4 保留的 benchmark 对象
\[
\mathcal Z_{sj}(M)=\{z\in Z:j\in\Gamma_s(z;M)\}
\]
在 ties 存在时可能重叠。Phase 7 没有把这个事实静默抹掉，而是明确写出：只有在可测 selector representation 的补充条件下，才可以得到互不相交的 selector regions 和行和为 1 的 transition shares。这样，后续阶段才能在不篡改 benchmark primitive definition 的前提下，使用一个良定义的 transition interface。

\subsection*{End-of-Section Recap}

\begin{itemize}
    \item What has been established mathematically: Phase 7 明确构造了测度值映射 $\Phi_M$，在完全写明的压缩条件下证明了不变测度存在且唯一，并给出了 selector-based transition shares 与 benchmark transition objects 的精确接口。
    \item What assumptions were used: Phase 3--5 继承下来的 Bellman 与 entry objects、Phase 6 的 market-clearing 环境，以及本阶段显式列出的 SD0--SD3 中的 selector 与算子压缩补充条件。
    \item What remains to be derived next: 下一阶段要把唯一的稳态对象、价格出清方程与进入条件合并成完整的 stationary competitive equilibrium 定义与存在性框架。
\end{itemize}
